% ===========================================================================
%  comune/title-user-guide.tex — title page and "Before you start"
% ===========================================================================
\begin{titlepage}
\thispagestyle{empty}
\begin{tikzpicture}[remember picture, overlay]

  \fill[lyraInchiostro] (current page.north west)
    rectangle ([yshift = -9mm]current page.north east);

  \node[anchor = north west, inner sep = 0pt]
    at ([shift = {(28mm, -42mm)}]current page.north west)
    {\begin{minipage}{125mm}
       {\sffamily\fontsize{10}{12}\selectfont
        \color{lyraAccentoScuro}\bfseries
        LYRA CAPITAL\par}
       \vspace{3mm}
       {\sffamily\fontsize{40}{44}\selectfont\color{lyraInchiostro}
        User\\guide\par}
       \vspace{5mm}
       {\color{lyraAccento}\rule{34mm}{1.6pt}\par}
       \vspace{6mm}
       {\sffamily\fontsize{12}{17}\selectfont\color{lyraInchiostro!62}
        What each screen does, in which order\\
        to use them, and what the numbers mean.\par}
     \end{minipage}};

  \node[anchor = south west, inner sep = 0pt]
    at ([shift = {(28mm, 34mm)}]current page.south west)
    {\begin{minipage}{125mm}
       \begin{tcolorbox}[enhanced, colback = lyraPannello,
                         colframe = lyraBordo, boxrule = 0.6pt, arc = 4pt,
                         left = 10pt, right = 10pt, top = 9pt, bottom = 9pt]
         {\sffamily\footnotesize\color{lyraInchiostro}%
          \textbf{How to read it.} Each chapter is one area of the
          application. Read it from the start, or go straight to the chapter
          for the screen in front of you.\par
          \vspace{4pt}
          \textbf{To change Lyra}, rather than use it, there is another
          document: the \emph{Guida dello sviluppatore} (developer guide, in
          Italian), in the same folder.\par}
       \end{tcolorbox}
     \end{minipage}};

  \node[anchor = south east, inner sep = 0pt, align = right]
    at ([shift = {(-24mm, 18mm)}]current page.south east)
    {\sffamily\scriptsize\color{lyraInchiostro!45}
     Document generated from \LaTeX\ source\\
     \today};

\end{tikzpicture}
\end{titlepage}

\cleardoublepage

\thispagestyle{empty}
\vspace*{18mm}

{\sffamily\Large\bfseries\color{lyraInchiostro}Before you start\par}
\vspace{3mm}
{\noindent\color{lyraAccento}\rule{\textwidth}{1pt}\par}
\vspace{7mm}

\noindent
Lyra is for building investment strategies, testing them on past data and
--- when they hold up --- running them on a real account. The three happen
in that order, and this guide follows the same order.

\vspace{4mm}
\noindent
In the text, \voce{This} is a menu item or a tab, \pulsante{This} a button to
press, \campo{this} a field to fill in, and \percorso{Settings \freccia Data}
means: open \voce{Settings}, then the \voce{Data} tab.

\vspace{6mm}
{\sffamily\large\bfseries\color{lyraInchiostro}The boxes\par}
\vspace{3mm}

\begin{nota}{Note}
A clarification. Skip it and the operation still works.
\end{nota}

\begin{pratica}{In practice}
The operation, step by step.
\end{pratica}

\begin{attenzione}{Warning}
Something can be lost here: data, money or time. Worth reading first.
\end{attenzione}

\vspace{6mm}
\begin{perche}{A word about the numbers}
No result shown by the application is a forecast. A backtest says how a
strategy \emph{would have} behaved on past data, with estimated costs; it
does not say how it will behave. The chapter on \voce{Streams} explains which
numbers to look at to tell a strategy that worked from one that got lucky.
\end{perche}

\cleardoublepage
